\subsection{}
Une horloge à pendule utilise les oscillations harmoniques simples d’un pendule pour compter le passage des secondes. Supposons que vous avez une horloge à pendule de grande précision. Vous déménagez sur la Lune, où le champ gravitationnel est beaucoup moins intense, et vous emmenez votre horloge avec vous. Quelle modification devriez-vous faire pour que l’horloge marque le temps aussi bien sur la Lune que sur Terre ?

\begin{enumerate}[label=\Alph*)]
    \item Rien : l’horloge fonctionnera aussi bien sans aucune modification.
    \item Augmenter la masse du pendule.
    \item Diminuer la masse du pendule.
    \item Allonger la corde du pendule.
    \item Raccourcir la corde du pendule.
\end{enumerate}
\noindent\textbf{Justification :} \hrulefill \\[2.5ex]
\rule{\linewidth}{0.4pt} \\[2.5ex]
\rule{\linewidth}{0.4pt}
\vfill